\clearpage

\section{相关理论基础}

\par 本章介绍本文工作所依赖的理论基础，包括差分隐私的形式化定义和常用机制、
图数据的基本概念和边差分隐私模型、Louvain 社区检测算法，以及作为本文工作直接基础的 PrivGraph 
算法的三阶段流程。

\subsection{差分隐私}

\subsubsection{差分隐私的形式化定义}

\par 差分隐私通过限制单条记录对算法输出的影响来提供隐私保护。
对于一个随机算法 $\mathcal{A}$，若对任意两个邻接数据集 $D$ 与 $D'$
（即 $D'$ 可由 $D$ 添加或删除一条记录得到，记为 $D \simeq D'$），
以及任意输出集合 $T \subseteq \mathrm{Range}(\mathcal{A})$，都有

\begin{equation}
    \label{equ:dp-def}
    \Pr[\mathcal{A}(D) \in T] \leq \exp(\varepsilon) \cdot \Pr[\mathcal{A}(D') \in T]
\end{equation}

\par 则称算法 $\mathcal{A}$ 满足 $\varepsilon$-差分隐私（$\varepsilon$-DP）。其中，参数 $\varepsilon$ 称为隐私预算（privacy budget），它控制着隐私保护的强度：$\varepsilon$ 越小，对相邻数据集输出分布的差异约束越严格，隐私保护越强；反之 $\varepsilon$ 越大，约束越松，可用性越好。差分隐私的本质保证是：攻击者即使获得了算法的完整输出和除某一条记录外的所有其他记录，其推断该记录是否在数据集中的概率比仍被 $e^{\varepsilon}$ 所限制，无法以显著优势加以区分。

\subsubsection{拉普拉斯机制}

\par 拉普拉斯机制（Laplace Mechanism, LM）是实现差分隐私最常用的机制之一，
适用于输出为实数或实数向量的查询函数。给定一个查询函数 $f : \mathcal{D} \to \mathbb{R}^k$，
定义其全局敏感度为

\begin{equation}
    \label{equ:lm-sens}
    \Delta_f = \max_{D \simeq D'} \| f(D) - f(D') \|_1
\end{equation}

\par 拉普拉斯机制对 $f$ 的输出的每个分量独立加上服从拉普拉斯
分布 $\mathrm{Lap}(\Delta_f / \varepsilon)$ 的噪声，
即 $\mathcal{A}(D) = f(D) + \mathrm{Lap}(\Delta_f / \varepsilon)$，
其中 $\mathrm{Lap}(b)$ 的概率密度函数为 $\frac{1}{2b}\exp\!\left(-\frac{|x|}{b}\right)$。可以证明，
拉普拉斯机制满足 $\varepsilon$-差分隐私。

\subsubsection{指数机制}

\par 指数机制（Exponential Mechanism, EM）适用于输出为离散值
（如类别、个体）的场景。给定一个数据集 $v$、候选输出集合 $\mathcal{O}$ 和
评分函数 $q : (v, o) \to \mathbb{R}$，定义评分函数的全局敏感度为

\begin{equation}
    \label{equ:em-sens}
    \Delta_q = \max_{o,\, v \simeq v'} | q(v, o) - q(v', o) |
\end{equation}

\par 指数机制以正比于 $\exp\left( \varepsilon \cdot q(v, o) / (2 \Delta_q) \right)$ 
的概率从 $\mathcal{O}$ 中采样输出 $o$。可以证明，指数机制满足 $\varepsilon$-差分隐私。

\subsubsection{差分隐私的组合性质}

\par 差分隐私具有良好的组合性质，使得复杂算法可以由简单的子机制组合而成。
常用的组合规则包括：

\begin{enumerate}
    \item 顺序组合（Sequential Composition）：如果算法 $\mathcal{A}_1, \mathcal{A}_2, \ldots, \mathcal{A}_k$ 分别满足 $\varepsilon_1, \varepsilon_2, \ldots, \varepsilon_k$-差分隐私，则它们的顺序组合满足 $(\varepsilon_1 + \varepsilon_2 + \cdots + \varepsilon_k)$-差分隐私。
    
    \item 并行组合（Parallel Composition）：如果算法 $\mathcal{A}_1, \mathcal{A}_2, \ldots, \mathcal{A}_k$ 分别作用于不相交的数据子集且分别满足 $\varepsilon_1, \varepsilon_2, \ldots, \varepsilon_k$-差分隐私，则它们的并行组合满足 $\max(\varepsilon_1, \varepsilon_2, \ldots, \varepsilon_k)$-差分隐私。
    
    \item 后处理性质（Post-processing）：对一个 $\varepsilon$-差分隐私算法的输出进行任意（不依赖于原始数据的）处理后得到的结果仍然满足 $\varepsilon$-差分隐私，即对任意 $g$ 满足 $g(\mathcal{A}(D))$ 也是 $\varepsilon$-差分隐私的。
\end{enumerate}

\par 后处理性质是本文工作的重要理论支撑：
本文对图重建阶段的所有改进均仅作用于已经噪声扰动过的中间结果（社区划分、噪声度序列、噪声边向量）
以及已经合成的图，不再访问任何原始数据。因此根据后处理性质，所有改进都不消耗任何额外的隐私预算，
整体算法的隐私性与原始 PrivGraph 完全一致。

\subsection{图数据与边差分隐私}

\subsubsection{图的基本概念}

\par 一张无向无权图可以记为 $G = (V, E)$，其中 $V$ 是节点集合，$E \subseteq V \times V$ 是
边集合。本文研究的图均为无向、无权、无自环的简单图。图可以用
邻接矩阵 $M \in \{0, 1\}^{|V| \times |V|}$ 表示，其中 $M[i][j] = 1$ 当且仅当存在一条连接
节点 $i$ 和节点 $j$ 的边。

\par 节点的度（degree）定义为该节点关联的边数。一张图的度序列即所有节点度数构成的向量。
度分布是图最基本的统计特征之一，许多真实图的度分布服从幂律（power law），即少数节点拥有非常多
的邻居（称为枢纽节点），大多数节点只有少量邻居。

\par 除度信息外，图还包含许多其他重要的结构特征，例如直径（diameter，所有节点对之间最短路径
的最大值）、聚类系数（clustering coefficient，反映节点周围邻居之间的连接紧密程度）、
模块度（modularity，反映社区划分质量）、特征向量中心性（eigenvector centrality, EVC，反映
节点的重要性）等。一个高质量的图合成方法应当在尽量多的图特征上保持原图的统计特性。

\subsubsection{边差分隐私}

\par 将差分隐私从一般数据集扩展到图数据上，关键在于定义两张图之间的``邻接关系''。
根据邻接关系的不同选择，差分隐私在图上有两种典型变体。

\par \textbf{边邻接（edge-neighboring）}\cite{nissim_smooth_2007}：
若两张图 $G = (V, E)$ 和 $G' = (V', E')$ 满足 $V = V'$ 且 $|E \oplus E'| = 1$，
即两者节点集相同、边集仅相差一条边，则称 $G$ 与 $G'$ 是边邻接的。
其中 $A \oplus B = (A \cup B) \setminus (A \cap B)$ 表示对称差。

\par \textbf{节点邻接（node-neighboring）}\cite{kasiviswanathan_analyzing_2013}：
若两张图中一张可以通过从另一张中移除一个节点及其所有关联边得到，
则称二者是节点邻接的。

\par 基于这两种邻接关系，差分隐私在图上分别衍生出节点差分隐私（node-DP）
与边差分隐私（edge-DP）：节点差分隐私提供更强的隐私保护，但全局敏感度通常较高，
难以获得高可用性的算法；边差分隐私提供相对适中的隐私保护，全局敏感度较低，
是当前实用方法采用的主流模型。

\par 形式化地，边差分隐私（$\varepsilon$-edge DP）定义为：
对于任意两张边邻接图 $G$ 和 $G'$，

\begin{equation}
    \label{equ:edge-dp}
    \Pr[\mathcal{A}(G) \in T] \leq \exp(\varepsilon) \cdot \Pr[\mathcal{A}(G') \in T], \quad \forall T \subseteq \mathrm{Range}(\mathcal{A})
\end{equation}

\par 本文与 PrivGraph 一致，均在 $\varepsilon$-边差分隐私模型下展开讨论。

\subsection{社区检测与 Louvain 算法}

\par 社区检测是发现图中稠密连通子结构的常用技术。直观上，一个社区指图中内部边稠密、
外部边稀疏的节点子集。社交网络中的兴趣群体、协作网络中的研究小组、生物网络中的功能模块都
可以用社区结构刻画。Newman 和 Girvan 提出的模块度（modularity）是衡量社区划分质量最常用的指标
，定义为

\begin{equation}
    \label{equ:modularity}
    Q = \sum_{C} \left[ \frac{\Sigma_{\mathrm{in}}(C)}{2m} - \left( \frac{\Sigma_{\mathrm{tot}}(C)}{2m} \right)^2 \right]
\end{equation}

\par 其中 $C$ 表示某个社区，$\Sigma_{\mathrm{in}}(C)$ 是社区 $C$ 内部所有边权重之和的两倍，$\Sigma_{\mathrm{tot}}(C)$ 是与社区 $C$ 内节点关联的所有边权重之和，$m$ 是图中所有边权重之和。模块度越大，说明社区划分越显著。

\par 在此基础上，引入分辨率参数 $t$ 后，模块度的扩展形式为：
\begin{equation}
    \label{equ:modularity-resolution}
    Q^{*} = \sum_{C} \left[ \frac{\Sigma_{\mathrm{in}}(C)}{2m} t - \left( \frac{\Sigma_{\mathrm{tot}}(C)}{2m} \right)^2 \right]
\end{equation}

\par 其中 $t$ 为分辨率参数，其他参数的含义与公式~\ref{equ:modularity} 一致。当 $t = 0$ 时，每个节点为一个独立社区；当 $t = 1$ 时，优化目标与传统模块度优化一致；随着 $t$ 的增大，社区数量通常会减少。
\par Louvain 算法是基于模块度优化的贪心层次聚类方法，由 Blondel 等人于 2008 年提出，因其算法简单、效果优良、扩展性好而被广泛使用。算法分为两个阶段反复迭代：

\par \textbf{阶段一（局部移动）：}对图中每个节点，计算将其移到每个邻居所在社区时的模块度增益，将该节点移动到使模块度增益最大的社区。如果不存在正增益，则保持不动。这一阶段反复扫描所有节点直至模块度收敛于局部最大值。

\par \textbf{阶段二（社区聚合）：}将每个社区视为一个超级节点，社区内部边的权重之和成为超级节点的自环权重，社区间的边权重成为超级节点之间的边权重。在新的超级图上重新执行阶段一。

\par Louvain 算法不断重复上述两个阶段直至模块度不再显著增加。然而，Louvain 算法的优化目标是模块度，存在著名的``分辨率限制''问题：在大型图中，模块度优化倾向于将小社区合并到大社区中，从而漏检小社区。该问题在 PrivGraph 中通过引入分辨率参数 $t$ 部分缓解。

\subsection{PrivGraph 算法}

\par PrivGraph 是Yuan等人在 USENIX Security 2023 上提出的差分隐私图发布方法，
是本文工作的直接基础。PrivGraph 的核心思想是利用图的社区结构作为信息抽取与噪声注入的中间粒度：
社区内部边稠密、社区间边稀疏，这一区分使得在不同粒度上分别处理两类边成为可能。PrivGraph 
的整体流程分为三个阶段：社区划分、信息抽取、图重建。
整体隐私预算 $\varepsilon = \varepsilon_1 + \varepsilon_2 + \varepsilon_3$ 
通常按 $1/3$ 平均分配。本文的优化工作集中在第三个阶段——图重建阶段。

\subsubsection{社区划分}

\par 社区划分阶段进一步分为社区初始化（Community Initialization）和社区调整（Community Adjustment）两步，
分别消耗 $\varepsilon_1$ 和 $\varepsilon_2$ 隐私预算。

\par 社区初始化的目标是获得一个粗粒度的初始社区划分。具体地，PrivGraph 
首先将所有节点随机划分为若干组，每组包含 $N$ 个节点。然后将每组节点视为一个超级节点，
由此构造一张加权图：每个超级节点有一个内权重，等于该组内所有节点度数之和；任意两个超级节点之间
有一条边，其外权重等于两组之间真实边的数量。PrivGraph 对内权重和外权重分别加拉普拉斯噪声以
满足 $\varepsilon_1$-边差分隐私（敏感度分别为 $2$ 和 $1$），并对噪声后的负值通过 NormSub 
后处理消除。最后，对噪声加权图运行 Louvain 算法，获得社区划分 $\mathbb{C}_W$，
并将其映射回原始节点空间得到 $\mathbb{C}_P$。

\par 社区调整以 $\mathbb{C}_P$ 为起点，对每个节点逐一应用指数机制：将其从原社区移除，
根据其与各候选社区的连边数量计算评分，按指数机制采样最终所属社区。该步骤消耗 $\varepsilon_2$ 
预算，得到最终的社区划分 $\mathbb{C}_F$。

\subsubsection{信息抽取}

\par \par 信息抽取阶段从原图中抽取两类信息：每个社区内部所有节点的度序列 $d$，
以及每对社区之间的边数构成的矩阵 $v$（即 $v(C_i, C_j)$ 表示社区 $C_i$ 与 $C_j$ 之间的边数）。这两类信息构成原图邻接矩阵的不相交子集，
因此可以并行使用同一份隐私预算 $\varepsilon_3$。具体地，
度序列的全局敏感度为 $2$（一条边变化影响两个节点的度），
边向量的敏感度为 $1$。PrivGraph 分别对 $d$ 和 $v$ 加拉普拉斯噪声并使用 NormSub 后处理，
得到噪声后的度序列 $\hat{d}$ 和边向量 $\hat{v}$。

\subsubsection{图重建}

\par 图重建阶段不再访问原图，仅基于噪声社区划分 $\mathbb{C}_F$、噪声度序列 $\hat{d}$ 和
噪声边向量 $\hat{v}$ 合成一张新图。该阶段属于差分隐私的后处理过程，不消耗任何隐私预算。
具体地，PrivGraph 采用 Chung-Lu（CL）模型分别处理两类边。

\par 对于社区内部边，PrivGraph 在每个社区 $C$ 内部根据 $\hat{d}$ 计算节点 $u$ 和
节点 $w$ 之间的连接概率

\begin{equation}
    \label{equ:cl-intra}
    p_{u, w} = \frac{\hat{d}_u \cdot \hat{d}_w}{\sum_{p \in C} \hat{d}_p}
\end{equation}

\par 并据此采样生成边。

\par 对于跨社区边，PrivGraph 仅有每对社区之间的边数 $\hat{v}(C_i, C_j)$ 作为约束，
丢失了端点的精确信息。因此 PrivGraph 采用均匀随机方式生成边：
直接在 $C_i$ 中均匀采样 $\hat{v}(C_i, C_j)$ 个起点，在 $C_j$ 中均匀采样相同数量的终点，
逐对连接。

\par 正是图重建阶段中跨社区边的均匀随机生成方式，
以及合成图在局部三角结构（聚类系数）和社区结构（模块度）上的偏差，
构成了本文优化的主要关注点。下一章将对此进行详细分析。

\clearpage

\section{PrivGraph 图重建阶段的局限性分析}
\label{chap:limitation}

\par 第二章回顾了 PrivGraph 的整体流程，可以看到 PrivGraph 在社区划分和信息抽取阶段的设计
相对精细：通过分级噪声注入、敏感度精准建模、NormSub 一致性后处理等手段，
将隐私预算的利用效率发挥到了较高的水平。然而，相比之下，
图重建阶段的设计则显得相对直接：社区内部边沿用 Chung-Lu 模型与噪声度序列，
跨社区边则简单地以均匀随机的方式连接节点；整个过程缺乏面向``局部结构''（如三角形）
和``全局结构''（如社区紧密度）的进一步约束。本章将对 PrivGraph 图重建阶段的不足进行系统性
的分析，从理论与实现两个层面剖析其问题成因，为下一章的优化设计提供动机。

\subsection{原始图重建算法回顾}

\par 为方便分析，本节首先以伪代码形式给出 PrivGraph 原始图重建阶段的核心实现。
该实现以噪声社区划分 $\mathbb{C}_F$、噪声度序列 $\hat{d}$、噪声边向量 $\hat{v}$ 为输入，
输出一张合成图的邻接矩阵 \texttt{mat2}。

\begin{algorithm}[H]
    \footnotesize 
    \caption{PrivGraph 原始图重建过程}
    \label{algo:step6-original}
    \begin{algorithmic}[1]
        \REQUIRE 空图 $G_N$，社区划分 $\mathbb{C}_F = \{C_1, \ldots, C_K\}$，
                 噪声度序列 $\hat{d}$，噪声边矩阵 $\hat{v}$
        \ENSURE 合成图 $G_s$
        \STATE 初始化 $G_s \gets G_N$
        \STATE \textbf{阶段 1：社区内部边生成}
        \FOR{$C_i \in \mathbb{C}_F$}
            \STATE 基于 $\hat{d}_i$ 生成 $C_i$ 内部子图 $s_i$ \quad // Chung-Lu 模型
            \STATE $G_s \gets G_s \cup s_i$
        \ENDFOR
        \STATE \textbf{阶段 2：跨社区边生成}
        \FOR{$C_i, C_j \in \mathbb{C}_F,\ i < j$}
            \STATE 在 $C_i \times C_j$ 上均匀随机采样 $\hat{v}_{ij}$ 条边构成 $s_{ij}$
            \STATE $G_s \gets G_s \cup s_{ij}$
        \ENDFOR
        \RETURN $G_s$
    \end{algorithmic}
\end{algorithm}

\par 算法~\ref{algo:step6-original} 给出了 PrivGraph 原始图重建过程的形式化描述。
其中输入为已发布的噪声社区划分 $\mathbb{C}_F = \{C_1, \ldots, C_K\}$、
噪声度序列 $\hat{d}$ 与噪声边向量 $\hat{v}$，
$\hat{d}_i$ 表示 $\hat{d}$ 在社区 $C_i$ 上的子序列，
$\hat{v}_{ij}$ 表示矩阵 $\hat{v}$ 中对应社区对 $(C_i, C_j)$ 的元素；
输出为合成图 $G_s$，初始化为节点数与原图相同的空图 $G_N$。
重建过程通过逐社区生成子图 $s_i$ 与 $s_{ij}$ 并合并入 $G_s$ 来逐步构造合成图。

\par 阶段 1 对每个社区 $C_i$，基于 Chung-Lu 模型~\cite{chung_connected_2002}
按噪声度序列 $\hat{d}_i$ 生成社区内部子图 $s_i$：每对节点 $(u, v) \in C_i \times C_i$
以概率 $\min(1, \hat{d}_u \hat{d}_v / (2 m_i))$ 独立连边，
其中 $m_i = \frac{1}{2}\sum_{p \in C_i} \hat{d}_p$。
阶段 2 对每对社区 $(C_i, C_j)$，从 $C_i$ 与 $C_j$ 中分别均匀独立采样
$\hat{v}_{ij}$ 个节点，逐对连接构成跨社区子图 $s_{ij}$。

\par 观察该算法可以发现，原始重建在两类边的处理上采用了两种不同的策略：
社区内部边基于 $\hat{d}$ 进行有偏采样，保留了节点度数的偏好性；
而跨社区边则采用均匀随机采样，整个过程未引入任何度信息。

\subsection{局限性分析}

\par 基于对原始算法的细致分析，本文识别出 PrivGraph 图重建阶段存在的三类典型问题：
跨社区边均匀采样导致的度信息利用不足、跨社区边对下游社区检测的过度影响，
以及合成图在局部三角结构和全局社区结构上的偏差。下面分别讨论。

\subsubsection{跨社区边均匀采样：度信息利用不足}
\label{subsubsec:cross-uniform}

\par PrivGraph 在信息抽取阶段同时记录了社区内度序列 $\hat{d}$ 和跨社区边数 $\hat{v}$，
但在重建跨社区边时仅使用了 $\hat{v}$ 这一标量信息，对每个端点都采用均匀随机采样。
这一设计的隐含假设是：在跨社区粒度上，节点之间的连接是无差异的。然而，
这一假设与真实图的统计特性存在明显偏离。

\par 真实图普遍存在所谓的``枢纽效应''或``偏好连接''现象：度数高的节点不仅在社区内部连接更密，
在跨社区方向上也更倾向于建立联系。从形式化的角度看，若记节点 $u$ 的度为 $d_u$，
则在 BA 模型、配置模型等经典图模型中，节点 $u$ 与节点 $v$ 形成边的概率均与 $d_u \cdot d_v$ 
成正比，这一规律在社区内和跨社区上是一致的。Chung-Lu 模型的连接概率也正是基于这一假设。

\par PrivGraph 在社区内部利用了这一规律（通过 \texttt{generate\_intra\_edge} 中
的度加权概率采样），但在跨社区方向上将其完全丢弃。这导致两个直接的后果：
第一，合成图中跨社区方向上的高度数节点（即枢纽节点）未能保持其``枢纽性''，
特征向量中心性等基于全局连接结构的指标会出现明显误差；第二，节点度数会均匀分布到
所有节点上，原图中度分布的长尾特性在跨社区方向上无法得到保留，进一步加剧了度分布 KL 散度的偏差。

\subsubsection{跨社区边对下游社区检测的过度影响}
\label{subsubsec:cross-overinfluence}

\par 原始 PrivGraph 算法在跨社区边生成阶段对每对社区 $(i, j)$ 按 
$\hat{v}(C_i, C_j)$ 给定的边数预算进行均匀采样，目的是在合成图中
还原原图的跨社区连接强度。然而可以观察到，跨社区边在 Louvain 等
模块度优化算法中具有特殊地位：它们直接影响算法在迭代过程中是否
将相邻社区合并的判断。在合成图上重新执行社区检测时，
过量的跨社区边会被算法判定为社区合并的信号，
使得新检测出的社区划分与原图差异较大。

\par 这一问题在差分隐私语境下被进一步放大。
原始 PrivGraph 中的社区划分本身已受指数机制噪声扰动而存在偏差，
$\hat{v}$也经过 Laplace 加噪，二者均非真值的精确反映。
若在此基础上仍然忠实地按 $\hat{v}$全量采样，
相当于在被噪声扰动过的社区划分上保留规模仍接近真值的跨社区连接，
两层误差叠加将进一步加剧合成图与原图在社区结构上的不一致。

\par 基于这一观察，本文将在下一章中提出一种``以适度牺牲跨社区边的统计还原度、
换取社区结构稳定性''的改进策略，以避免噪声桥梁对下游社区检测的过度干扰。

\subsubsection{合成图局部三角结构的固有局限}

\par 除上述两个问题外，PrivGraph 在生成社区内部边时基于 Chung-Lu 模型，
仅利用了节点度数信息，并未显式建模"邻居之间的关联"。
这就导致合成图中虽然每个节点的度数与目标 $\hat{d}$ 基本一致，
但边的位置是按度概率随机分布的，缺乏真实图所特有的传递性。
形式化地说，真实社交网络中"朋友的朋友也是朋友"这一现象使得节点的
聚类系数普遍较高，而 Chung-Lu 类的随机图模型理论上聚类系数趋于 $0$，
这是由模型本身决定的偏差。

\par 这一问题与跨社区边过多的问题相互耦合：
跨社区边均匀分布在每对社区之间，使得社区与社区之间的"边界"在合成图中
显得模糊，整体模块度低于原图，进一步削弱了合成图的局部传递性表现。
不过需要指出的是，三角结构的根本性修复需要超出 Chung-Lu 框架的额外建模手段，
因此本文将这一问题作为一个开放方向，
而将主要优化精力集中于跨社区边处理这一更为关键的瓶颈环节。

\subsubsection{局限性的总体影响}

\par 三类问题之间并不是孤立的，而是存在协同放大关系。跨社区均匀采样使得跨社区边
在节点上均匀分布，未能反映原图中跨社区连接的度异质性；
这种偏差又被原图重建过程中已存在的双层噪声（噪声社区划分与噪声边向量 $\hat{v}$）
所固化，最终在 NMI、模块度等指标上同时表现出系统性偏差。
其中，跨社区边的处理同时关联前两个问题——
既涉及度信息利用，又直接影响下游社区检测——
是图重建阶段最为关键的瓶颈环节。本文据此将优化重点集中于跨社区边的处理，
通过单一但针对性较强的图重建模块同时缓解前两类问题；
对于第三类局部三角结构偏差，本文将其作为一个开放方向，留作后续研究。

\clearpage
\section{跨社区边处理的图重建优化}

\par 本章详细介绍针对第三章所识别的局限性所提出的图重建优化方法。
\autoref{subsec:design-overview}介绍整体设计思路；
\autoref{subsec:reconstruct}介绍图重建模块的具体设计，
针对原 PrivGraph 在跨社区边均匀采样过程中度信息利用不足、
以及与噪声社区划分耦合放大的问题进行重新设计；
\autoref{subsec:analysis}给出隐私性证明与复杂度分析。

\subsection{设计思路与总体框架}
\label{subsec:design-overview}

\par 本文优化的总体设计遵循以下三条思路：

\begin{enumerate}
    \item \textbf{后处理性质的利用}\cite{dwork_algorithmic_2014}\textbf{：}
    所有优化均仅作用于差分隐私的后处理阶段，
    即只能访问已经被噪声扰动过的中间结果（$\mathbb{C}_F$、$\hat{d}$、$\hat{v}$）
    以及已经合成的图本身，不再访问任何原始数据。
    根据差分隐私的后处理性质（Post-processing Property），
    对 $\varepsilon$-DP 算法输出做不依赖于原始数据的任意处理，
    结果仍满足 $\varepsilon$-DP。这一性质保证本文优化算法的隐私性与原始 PrivGraph 
    完全一致，因此无需重新进行隐私预算分析。
    
    \item \textbf{与原算法的模块化兼容性：}本文优化与原始 PrivGraph 在结构上保持一致，
    仍然以``社区内部边 + 跨社区边''为主要划分，
    新增的图重建模块以独立函数封装并提供显式开关，
    便于作为原算法的``可替换组件''集成到已有 PrivGraph 实现中。
    
    \item \textbf{以瓶颈环节为优化对象：}\autoref{chap:limitation} 的局限性分析表明，
    PrivGraph 图重建阶段中跨社区边的处理是合成图可用性的主要瓶颈，
    对 NMI、模块度等多项指标存在系统性影响。
    本文据此将优化重点放在跨社区边的处理上，
    构成 \autoref{subsec:reconstruct} 节所述的图重建模块。
\end{enumerate}

\par 基于上述原则，本文针对图重建阶段的跨社区边处理进行了重新设计：
原算法以噪声边向量 $\hat{v}$ 为预算上限、对每对社区独立执行均匀采样，
既未利用 $\hat{d}$ 中蕴含的度信息，也未对采样规模与噪声社区划分之间的
耦合误差加以控制。本文在该模块中引入两项相互配合的改动：
将跨社区采样从均匀切换为基于 $\hat{d}$ 的度加权有偏采样，
使噪声度序列在跨社区方向上同样得到利用；
按比例 \texttt{inter\_ratio} 主动缩减每对社区之间生成的跨社区边数，
削弱噪声社区划分与噪声 $\hat{v}$ 两层误差对合成图社区结构的叠加影响。
此外，为补偿因跨社区边缩减而引起的总边数下降，
本文在 Chung-Lu 采样之前根据每个社区节省的跨社区边数
对噪声度序列 $\hat{d}$ 进行等比例放大，
使 Chung-Lu 模型在社区内部自然多生成相应数量的边，
从而在保持度分布形状不变的前提下维持节点度分布的整体一致性。

\par 第五章将通过实验对该模块的效果进行系统评估。

\subsection{图重建模块的优化}
\label{subsec:reconstruct}

\subsubsection{设计动机}

\par \autoref{subsubsec:cross-uniform} 与 \autoref{subsubsec:cross-overinfluence} 指出，PrivGraph 的图重建阶段在跨社区边处理上
存在两个相互关联的问题：跨社区采样未利用噪声度序列 $\hat{d}$，
导致跨社区方向上度信息丢失；而跨社区边数完全由噪声边向量 $\hat{v}$ 给定，
在噪声社区划分本身已存在偏差的前提下，忠实按 $\hat{v}$ 全量采样会使两层噪声叠加，
进而使合成图的社区结构在评测时表现出系统性偏差。这两类问题在低 $\varepsilon$ 区段
被同时放大——$\hat{v}$ 本身方差较大，叠加噪声社区划分的边界偏差后，
合成图与原图在社区结构上的差异被显著抬升。

\par 与之相伴的是社区内部的结构稳定性问题。原算法直接
按 $\hat{d}$ 进行 Chung-Lu 采样，该过程未显式控制采样涨落造成的边比例偏差；
而 Chung-Lu 模型本身缺乏对三角结构的建模，使社区内部的实际边分布与原图存在系统性偏离。
跨社区与社区内部两类问题的耦合，使合成图在度分布、聚类系数与模块度等多项指标上
同时表现出可观测的误差。

\par 由于跨社区边的处理同时影响整体连通性、模块度、NMI 等多个关键指标，
本文将其作为图重建模块的主要优化对象。该模块的基本思路是：
将原 PrivGraph 中均匀的跨社区采样替换为基于 $\hat{d}$ 的度加权有偏采样，
使噪声度序列在跨社区方向上同样得到利用；
以 \texttt{inter\_ratio} 主动缩减每对社区之间的跨社区边规模，
削弱``噪声社区划分 $\times$ 噪声 $\hat{v}$''的双层误差叠加效应；
同时，在 Chung-Lu 采样之前根据每个社区因跨社区缩减而节省的边数
对 $\hat{d}$ 进行等比例放大，使 Chung-Lu 模型在社区内部自然补偿
因跨社区边减少而损失的总边数，在保持度分布形状不变的前提下
维持节点度分布的整体一致性。
\par 第五章的主对比实验将进一步验证：本模块在 NMI、Modularity 等指标上
能够稳定改善合成图的可用性。

\subsubsection{算法流程}

\par 图重建模块的整体过程分为三个阶段。第一阶段根据跨社区边的缩减量
对噪声度序列进行等比例放大；第二阶段使用放大后的度序列进行 Chung-Lu 社区内部边生成；
第三阶段对跨社区边进行度加权缩减采样。

\begin{enumerate}
    \item \textbf{度序列放大（边预算补偿）：}对每个社区 $C_i$，
    计算其因跨社区缩减而节省的边数。具体地，社区 $C_i$ 与所有其他社区之间
    原本需要生成的跨社区边数为 $\sum_{j \neq i} \hat{v}(C_i, C_j)$，
    按 \texttt{inter\_ratio} 缩减后节省的边数为
    $\mathrm{extra}_i = \sum_{j \neq i} \hat{v}(C_i, C_j) \cdot (1 - \texttt{inter\_ratio}) / 2$
    （除以 $2$ 是因为每对社区的节省量在两个社区之间均分）。
    由于 Chung-Lu 模型生成的期望边数约为 $\sum_{p \in C_i} \hat{d}_p / 2$，
    要使社区内多生成 $\mathrm{extra}_i$ 条边，需要将度序列之和增大 $2 \cdot \mathrm{extra}_i$。
    因此对 $C_i$ 中每个节点的度数按统一的放大系数
    \begin{equation}
        \label{equ:scale-factor}
        \mathrm{scale}_i = 1 + \frac{2 \cdot \mathrm{extra}_i}{\sum_{p \in C_i} \hat{d}_p}
    \end{equation}
    进行等比例放大，并将放大后的度数截断至 $[0,\, |C_i| - 1]$ 以保持合法性。
    该操作保持了度分布的形状不变，仅整体提升了边的生成量。

    \item \textbf{社区内部边生成：}对每个社区 $C_i$，
    基于放大后的度序列 $\hat{d}_i^{\mathrm{scaled}}$ 调用 Chung-Lu 采样
    生成社区内部边，置入合成邻接矩阵 \texttt{mat2} 的相应块。
    由于度序列已经包含了边预算补偿，Chung-Lu 模型在社区内部自然多生成
    与跨社区缩减量相当的边数，无需额外的随机补边操作。

    \item \textbf{跨社区度加权缩减采样：}对每对社区 $(i, j)$，
    按 \texttt{inter\_ratio} 计算保留的跨社区边数 
    $n_{\mathrm{inter}} = \lfloor \hat{v}(C_i, C_j) \cdot \texttt{inter\_ratio} \rfloor$。
    从 $\hat{d}_i$、$\hat{d}_j$ 归一化得到概率向量 $\pi_i$、$\pi_j$，
    在 $C_i \times C_j$ 上按 $(\pi_i, \pi_j)$ 独立采样节点对，
    通过碰撞检测避免重复边与已存在边，
    直至成功填入 $n_{\mathrm{inter}}$ 条边。
    与原算法直接采样 $\hat{v}(C_i, C_j)$ 条均匀跨社区边相比，
    这一步在两方面对原流程进行了改造：
    一是将跨社区采样从均匀切换为度加权，
    使 $\hat{d}$ 在跨社区方向上同样得到利用；
    二是按 \texttt{inter\_ratio} 主动缩减跨社区边规模，
    以缓解噪声社区划分与噪声 $\hat{v}$ 双层误差的叠加扰动。
\end{enumerate}

\par 其中 \texttt{inter\_ratio} 是该模块唯一的核心超参，
直接决定跨社区边的保留比例，同时也间接决定了度序列放大的幅度。
本文取 $\texttt{inter\_ratio} = 0.10$，
即每对社区之间仅保留约 $\hat{v}(C_i, C_j) \cdot 10\%$ 条度加权采样的跨社区边，
其余 $90\%$ 的边预算通过度序列放大回流至社区内部的 Chung-Lu 采样。

\par 综合上述三个阶段，本文图重建优化模块的完整流程如算法~\ref{algo:reconstruct} 所示。

\begin{algorithm}[H]
    \footnotesize 
    \caption{图重建模块优化}
    \label{algo:reconstruct}
    \begin{algorithmic}[1]
        \REQUIRE 空图 $G_N$,社区划分 $\mathbb{C}_F = \{C_1, \ldots, C_K\}$,
                 噪声度序列 $\hat{d}$,噪声边矩阵 $\hat{v}$,
                 跨社区保留比例 $\eta$(即 \texttt{inter\_ratio})
        \ENSURE 合成图 $G_s$
        \STATE 初始化 $G_s \gets G_N$
        \STATE \textbf{阶段 0:度序列放大(边预算补偿)}
        \FOR{$C_i \in \mathbb{C}_F$}
            \STATE 由 $\hat{v}_{ij}$ 与 $\eta$ 按公式~(\ref{equ:scale-factor}) 
            计算 $\mathrm{scale}_i$,得 $\hat{d}_i^{\mathrm{scaled}} \gets \mathrm{scale}_i \cdot \hat{d}_i$
            (截断至 $[0,|C_i|-1]$)
        \ENDFOR
        \STATE \textbf{阶段 1:社区内部边生成}
        \FOR{$C_i \in \mathbb{C}_F$}
            \STATE 基于 $\hat{d}_i^{\mathrm{scaled}}$ 用 Chung-Lu 模型生成 $s_i$,
                   $G_s \gets G_s \cup s_i$
        \ENDFOR
        \STATE \textbf{阶段 2:跨社区度加权缩减采样}
        \FOR{$C_i, C_j \in \mathbb{C}_F,\ i < j$}
            \STATE 取 $n_{ij} \gets \lfloor \eta \cdot \hat{v}_{ij} \rfloor$,
                   按 $(\pi_i, \pi_j) \propto (\hat{d}_i, \hat{d}_j)$ 在 $C_i \times C_j$ 
                   上采样 $n_{ij}$ 条边构成 $s_{ij}$,
                   $G_s \gets G_s \cup s_{ij}$
        \ENDFOR
        \RETURN $G_s$
    \end{algorithmic}
\end{algorithm}

\subsection{整体流程、隐私性与复杂度}
\label{subsec:analysis}

\subsubsection{隐私性分析}
\par \textbf{命题 4.1}：如果 PrivGraph 的社区划分阶段
（消耗 $\varepsilon_1 + \varepsilon_2$ 隐私预算）和信息抽取阶段
（消耗 $\varepsilon_3$ 隐私预算）满足 
$(\varepsilon_1 + \varepsilon_2 + \varepsilon_3)$-边差分隐私，
则将其图重建阶段替换为本文优化模块后，
整个流程仍然满足 $(\varepsilon_1 + \varepsilon_2 + \varepsilon_3)$-边差分隐私。

\par \textbf{证明}：根据差分隐私的后处理性质，
对任何 $\varepsilon$-DP 算法的输出做不依赖原始数据的处理，
结果仍满足 $\varepsilon$-DP。下面验证图重建模块为纯后处理。

\par 图重建模块的输入为 
$(\mathbb{C}_F, \hat{d}, \hat{v}, \eta)$。
其中社区划分 $\mathbb{C}_F = \{C_1, \ldots, C_K\}$ 已满足 $(\varepsilon_1 + \varepsilon_2)$-DP；
噪声度序列 $\hat{d}$ 与噪声边向量 $\hat{v}$ 均已满足 $\varepsilon_3$-DP；
跨社区保留比例 $\eta$ 是公开超参；
节点数 $N = |\bigcup_i C_i|$、社区数 $K = |\mathbb{C}_F|$ 均由 $\mathbb{C}_F$ 派生（公开常数）。
模块内部派生的放大度序列 $\hat{d}^{\mathrm{scaled}}$ 仅由 $\hat{d}$、$\hat{v}$ 与 $\eta$ 计算得到，
整个重建过程在执行中未访问原图 $G$，因此该模块属于纯后处理。

\par 根据后处理性质，本文图重建模块的输出仍满足
$(\varepsilon_1 + \varepsilon_2 + \varepsilon_3)$-边差分隐私。
结合前置阶段，整体算法满足 $(\varepsilon_1 + \varepsilon_2 + \varepsilon_3)$-边差分隐私，
与原始 PrivGraph 完全一致。证毕。

\subsubsection{复杂度分析}

\par \textbf{时间复杂度}方面：
图重建模块新增的度序列放大操作对每个社区遍历一次度序列，代价为 $O(N)$，
可以忽略；Chung-Lu 采样与跨社区采样的主体部分中，
跨社区采样虽从均匀采样切换为度加权采样，但单次代价不变，
且采样总量由 $\eta$ 缩减，
因此该模块整体时间复杂度仍为 $O(N^2)$，与原算法相同。

\par \textbf{空间复杂度}方面：
原 PrivGraph 论文给出的理论空间复杂度为 $O(m + N)$，
但其官方实现采用稠密邻接矩阵存储合成图，实际空间为 $O(N^2)$。
本文沿用相同的稠密存储方式以保持实现的可比性。
图重建模块不引入额外存储，仅在合成图邻接矩阵 $A$ 上原地操作，
整体峰值仍为 $O(N^2)$，与原算法实现同阶。

\par 综合来看，本文优化算法在保留原始 PrivGraph 隐私性与算法复杂度阶的前提下，
通过图重建模块的优化改善了合成图的可用性。
该模块的贡献将在第五章~\autoref{subsec:exp-main-cha} 的主对比实验中得到验证。

\par 至此，本文优化方法的设计、隐私性与复杂度均已论述完毕。
下一章将通过实验对其有效性进行系统评估。

\clearpage
\section{实验评估}
\label{chap:experiment}

\par 本章对前文提出的优化方法进行系统的实验评估。
\autoref{subsec:exp-setup}介绍实验设置，包括数据集、对比方法、评测指标与超参配置；
\autoref{subsec:exp-main}给出本文方法与原始 PrivGraph 在多个数据集上的整体对比；
\autoref{subsec:exp-hyper}给出对核心超参 \texttt{inter\_ratio} 的敏感性分析；
\autoref{subsec:exp-summary}对实验结果进行小结。

\subsection{实验设置}
\label{subsec:exp-setup}

\subsubsection{数据集}

\par 本文在 Chameleon、Facebook、CA-HepPh、Enron 四个真实图数据集上展开实验。
四者覆盖了不同来源（社交网络、合作网络、邮件网络）、不同规模与不同稀疏程度的图，
是差分隐私图发布领域的常用基准。
其中 Chameleon 规模适中、社区结构鲜明，
便于对核心方法和超参进行细粒度对比，
本文以 Chameleon 为主数据集进行完整指标对比与超参敏感性分析；
其余三个数据集用于验证方法在不同图上的可推广性。
四个数据集的基本统计信息如\autoref{tab:dataset}所示。

\begin{table}[H]
    \caption{\label{tab:dataset}实验数据集的基本统计信息}
    \centering
    \begin{tabularx}{\linewidth}{|>{\centering\arraybackslash}X|>{\centering\arraybackslash}X|>{\centering\arraybackslash}X|>{\centering\arraybackslash}X|}
        \hline
        \textbf{数据集} & \textbf{节点数} & \textbf{边数} & \textbf{平均度} \\ \hline
        Chameleon & 2{,}277  & 31{,}371  & 27.55 \\ \hline
        Facebook  & 4{,}039  & 88{,}234  & 43.69 \\ \hline
        CA-HepPh  & 12{,}008 & 118{,}489 & 19.74 \\ \hline
        Enron     & 33{,}696 & 180{,}811 & 10.73 \\ \hline
    \end{tabularx}
\end{table}

\subsubsection{对比方法}

\par 本文将以下两种方法纳入对比：

\begin{enumerate}
    \item \textbf{PrivGraph}：原始算法，作为基线方法。
    其图重建阶段采用社区内 Chung-Lu 采样 + 跨社区均匀采样的组合，
    隐私预算按 $1{:}1{:}1$ 平均分配。
    
    \item \textbf{Ours}：本文方法。
    在 Chung-Lu 采样前根据跨社区缩减量对 $\hat{d}$ 进行等比例放大，
    跨社区采样改为基于 $\hat{d}$ 的度加权有偏采样并按 \texttt{inter\_ratio} 缩减。
    其余阶段（社区划分、信息抽取）保持与 PrivGraph 一致。
\end{enumerate}

\par 两种方法的隐私预算分配方式保持一致（$\varepsilon_1 = \varepsilon_2 = \varepsilon_3 = \varepsilon/3$），
以确保可比性。

\subsubsection{评测指标}

\par 本文沿用 PrivGraph 原文的评测体系，从社区结构、节点重要性、
度分布、全局结构与局部结构五个维度共选取 7 个指标：

\begin{itemize}
    \item \textbf{NMI}（归一化互信息）：合成图与原图在社区划分上的一致性，越高越好。
    \item \textbf{Modularity Relative}（模块度相对误差）：定义为 $|Q_{\mathrm{syn}} - Q_{\mathrm{ori}}| / Q_{\mathrm{ori}}$，越低越好。
    \item \textbf{EVC Overlap}：原图与合成图按特征向量中心性排序的前 $1\%$ 节点的重合比例，越高越好。
    \item \textbf{EVC MAE}：原图与合成图按特征向量中心性排序的前 $1\%$ 节点上的平均绝对误差，越低越好。
    \item \textbf{Degree KL}（度分布 KL 散度）：合成图与原图度分布之间的 KL 散度，越低越好。
    \item \textbf{Diameter Relative}（直径相对误差）：定义为 $|D_{\mathrm{syn}} - D_{\mathrm{ori}}| / D_{\mathrm{ori}}$，越低越好。
    \item \textbf{CC Relative}（聚类系数相对误差）：定义为 $|C_{\mathrm{syn}} - C_{\mathrm{ori}}| / C_{\mathrm{ori}}$，越低越好。
\end{itemize}

\subsubsection{超参与运行配置}

\par 实验在 7 个隐私预算 $\varepsilon \in \{0.5, 1.0, 1.5, 2.0, 2.5, 3.0, 3.5\}$
上展开。社区初始化的分组规模 $N = 20$，分辨率参数 $t = 1.0$，
均沿用 PrivGraph 默认设置。
本文方法的唯一核心超参取值为 $\texttt{inter\_ratio} = 0.10$。
主对比与可推广性实验对每组 $\varepsilon$ 独立运行 10 次后取均值；
\autoref{subsec:exp-hyper} 的超参敏感性扫描中
\texttt{inter\_ratio} 的每个取值独立运行 30 次，
以便对相近取值之间的差异进行 Welch t 检验。
所有实验均在同一台机器上以单线程方式完成。

\subsection{主对比实验}
\label{subsec:exp-main}

\par 本节给出本文方法 Ours 与原始 PrivGraph 在多个数据集上的整体对比。
\autoref{subsec:exp-main-cha}给出 Chameleon 数据集上两种方法在 7 个指标 × 7 个 $\varepsilon$
上的完整对比；\autoref{subsec:exp-main-other}在 Facebook、CA-HepPh 与 Enron 数据集上
对比两种方法在 7 个 $\varepsilon$ 下的核心指标，验证本文方法的可推广性。

\subsubsection{Chameleon 数据集上的完整对比}
\label{subsec:exp-main-cha}

\par 由 \autoref{fig:cha-overview} 可见，本文方法 Ours 在 Chameleon 数据集上
相对原始 PrivGraph 取得了显著改善，具体表现按指标维度分述如下：

\begin{figure}[H]
    \centering
    \includegraphics[width=\linewidth]{chamelon_overview7.pdf}
    \caption{\label{fig:cha-overview}PrivGraph 与本文方法在 Chameleon 数据集上 7 个指标的对比}
\end{figure}

\begin{itemize}
    \item \textbf{社区结构指标：}
    NMI 上 Ours 在所有 7 个 $\varepsilon$ 上均严格领先 PrivGraph，
    平均值由 0.160 提升至 0.309（相对提升约 93\%）；
    Modularity Relative 平均值由 0.395 降至 0.121（相对降低约 69\%），
    全段稳定领先。这一显著改善印证了 \autoref{subsec:reconstruct} 所述的设计动机——
    将均匀采样替换为基于 $\hat{d}$ 的度加权采样、并按 \texttt{inter\_ratio} 主动缩减
    跨社区边规模，有效削弱了``噪声社区划分 $\times$ 噪声 $\hat{v}$''的双层误差叠加效应。

    \item \textbf{度分布指标：}
    Degree KL 平均值由 4.12 降至 1.56（相对降低约 62\%），
    在低、中预算区段改善尤为显著，
    印证了将 $\hat{d}$ 用于跨社区方向、并通过度序列放大进行边预算补偿的有效性；
    高 $\varepsilon$ 区段两者差距缩小，符合预期——
    当 $\hat{d}$ 噪声较小时，原 PrivGraph 已能较好恢复度分布。

    \item \textbf{全局结构指标：}
    Diameter Relative 平均值由 0.518 降至 0.390（相对降低约 25\%），
    全段稳定领先，说明对跨社区边规模的合理控制对全局可达结构的还原同样有正向贡献。

    \item \textbf{节点重要性与局部结构指标：}
    EVC Overlap 与 EVC MAE 上两者表现整体相当
    （EVC Overlap 0.757 vs 0.760、EVC MAE $5.4 \times 10^{-3}$ vs $4.6 \times 10^{-3}$，
    差异在随机方差量级内）。
    CC Relative 上 Ours 平均值 0.279 略高于 PrivGraph 的 0.206，差距约 0.07。
    一种合理的解释是：本文方法的度序列放大策略在补偿总边数的同时，
    使社区内部 Chung-Lu 采样的边数也按比例增加，
    合成图的聚类系数与原图的偏离反而轻微增大。
    Chung-Lu 模型本身缺乏对三角结构的显式建模，
    这一现象与 \autoref{chap:limitation} 中关于局部三角结构作为开放问题的讨论一致。
\end{itemize}

\par 综合来看，本文方法在涉及社区结构、度分布与全局结构的核心指标上稳定占优，
在节点重要性与局部传递性指标上保持中性或略优，
与 \autoref{subsec:design-overview} 所述的设计目标一致。


\subsubsection{其他数据集上的可推广性验证}
\label{subsec:exp-main-other}

\par 为验证本文方法在不同规模与不同来源图上的可推广性，
本节在 Facebook、CA-HepPh 与 Enron 三个数据集上对 PrivGraph 与 Ours
进行 7 个 $\varepsilon$ 上的对比。考虑篇幅，
仅报告 NMI、Modularity Relative、Degree KL 与 CC Relative 四个最具代表性的指标，
分别对应社区结构、度分布与局部传递性三个维度。
结果如 \autoref{fig:multi-fbcah} 与 \autoref{fig:multi-enron} 所示。

\begin{figure}[H]
    \centering
    \includegraphics[width=\linewidth]{facebook_overview7.pdf}
    \includegraphics[width=\linewidth]{cahepph_overview7.pdf}
    \caption{\label{fig:multi-fbcah}Facebook 与 CA-HepPh 数据集上 PrivGraph 与本文方法 Ours 在七个核心指标上的对比}
\end{figure}

\par \textbf{多数据集结果分析：}
对照 \autoref{fig:multi-fbcah} 与 \autoref{fig:multi-enron}，
本文方法在三个不同来源、不同规模的数据集上
均能稳定改善社区结构相关指标，呈现一致的可推广性。
在规模适中、社区结构鲜明的 Facebook 上，
Ours 在 NMI、Modularity Relative、Degree KL、CC Relative 四项指标上均稳定优于 PrivGraph：
NMI 平均由 0.153 提升至 0.275（相对提升约 79\%），
Modularity Relative 平均由 0.485 降至 0.209（相对降低约 57\%），
Degree KL 平均由 3.98 降至 0.63（相对降低约 84\%），
CC Relative 平均由 0.622 降至 0.301（相对降低约 52\%），
表现稳定且趋势一致，是本文方法最具代表性的成功案例。
在合作网络 CA-HepPh 上，本文方法同样取得了全段改善：
NMI 平均由 0.066 提升至 0.376（5 倍以上），
Modularity Relative 平均由 0.384 降至 0.157（相对降低约 59\%），
Degree KL 平均由 3.42 降至 1.36（相对降低约 60\%），
CC Relative 平均由 0.555 降至 0.334（相对降低约 40\%）。
其中 NMI 的相对优势尤为悬殊，这与 CA-HepPh 强社区结构、
原 PrivGraph 又因跨社区均匀采样使社区边界均匀分布的特性吻合：
当原图社区结构越清晰、跨社区均匀采样产生的扰动越大时，
本文方法基于度加权采样并主动缩减跨社区规模的优势就越突出。

\begin{figure}[H]
    \centering
    \includegraphics[width=\linewidth]{enron_overview7.pdf}
    \caption{\label{fig:multi-enron}Enron 数据集上 PrivGraph 与本文方法 Ours 在七个核心指标上的对比}
\end{figure}

\par 在规模更大的邮件通信网络 Enron 上，本文方法整体仍保持显著优势，
但在 CC Relative 一项上呈现出有别于前两个数据集的反常表现。
社区结构方面，PrivGraph 在 7 个 $\varepsilon$ 上的 NMI 均值仅为 0.017，
几乎无法恢复原图的社区结构；本文方法则稳定维持在 0.33--0.35 之间，
相对提升超过一个量级，是三个数据集中改善幅度最大者。
Modularity Relative 由 PrivGraph 的均值 0.435 降至本文方法的 0.297，
相对降低约 32\%，其中 $\varepsilon \geq 2.5$ 时尤为明显——
PrivGraph 在此区间始终保持在 0.38--0.40 附近，而本文方法已下降至 0.14--0.22。
Degree KL 在低隐私预算下本文方法显著占优（$\varepsilon=0.5$ 时为 2.04 vs.\ 3.38），
但随 $\varepsilon$ 增大，PrivGraph 因未额外引入度信息扰动而逐渐反超，
约在 $\varepsilon=2.0$ 处出现交叉，
反映出度加权采样在噪声幅度较小时对噪声度序列的扰动更为敏感。

\par 与上述指标形成鲜明对照的是 CC Relative：
本文方法在 7 个隐私预算下的均值为 2.93，相较 PrivGraph 的 0.44 显著偏高，
尤其在 $\varepsilon \leq 1.0$ 区间达到 5--6 的量级。
究其原因，可归结为以下三方面：
其一，Enron 作为真实邮件通信网络，呈现典型的 hub-and-spoke 拓扑特征，
少数核心节点度数极高，且原图整体聚类系数较高；
其二，本文方法的两项核心机制——基于噪声度序列的度加权跨社区采样
与社区内噪声度序列的等比例放大补偿——共同作用于高度数节点，
使合成图中以 hub 为中心的星形子结构被进一步强化，
而星形子图的局部聚类系数趋近于 0，
从而系统性地拉低了合成图的整体 CC；
其三，CC Relative 以原图 CC 为分母进行归一化，
Enron 较高的原始 CC 进一步放大了上述偏差在相对指标上的呈现。
相比之下，PrivGraph 的均匀随机采样虽在社区结构指标上表现欠佳，
但其在各节点间分配边数较为均衡，反使合成图的 CC 与原图差距较小。

\par 一种可行的缓解思路是调整 PrivGraph 在``社区划分—信息抽取—图重建''
三个阶段间的隐私预算分配比例。作为初步验证，
本文将默认的均等分配 $(\varepsilon_1, \varepsilon_2, \varepsilon_3) = (1/3, 1/3, 1/3)$
调整为偏重信息抽取阶段的 $(0.25, 0.50, 0.25)$，
即将更多预算分配于度序列与边数抽取，以降低噪声度序列的扰动幅度。
Enron 上的实验结果表明，该调整可在多个隐私预算下稳定降低 CC Relative：
$\varepsilon=2.0$ 时由 2.42 降至 1.68，相对降幅约 31\%；
$\varepsilon=3.0$ 时由 1.50 降至 0.96，相对降幅约 36\%；
与此同时，NMI、Modularity Relative 等核心指标基本不受影响，部分预算下甚至略有提升。
上述结果表明，CC Relative 偏高的问题在相当程度上源于噪声度序列扰动
经度加权采样与等比例放大机制后的累积效应，
而将更多隐私预算分配于度序列抽取阶段即可有效抑制该效应。
三阶段隐私预算的联合最优分配本身构成一个独立的优化问题，
将作为后续工作进一步探索。

\par 综合三个数据集的结果可见，本文方法在不同规模、不同来源、
不同度分布特性的图上均能稳定可推广，
核心优势——基于度信息的跨社区采样与边预算补偿——在多种拓扑结构上都有效；
而对于 Enron 等具有显著 hub 结构、高聚类特性的网络，
通过适配性地调整三阶段隐私预算分配比例，
还可在保持社区结构相关指标既有优势的同时，
进一步抑制聚类系数指标上的副作用。





\subsection{超参敏感性分析}
\label{subsec:exp-hyper}


\par 本文方法引入唯一核心超参 \texttt{inter\_ratio}，
其取值决定每对社区上跨社区边的保留比例，是影响合成图整体连通性与社区结构的
关键变量。本节针对该超参在其默认值附近进行一维网格扫描，评估其敏感性。
扫描以 $\varepsilon = 2.0$ 为代表预算,对每个取值独立运行 30 次，
便于通过 Welch t 检验对相近取值之间的差异进行显著性评估。
本文在 $\{0.05, 0.10, 0.15, 0.20, 0.30\}$ 上进行扫描，
结果如\autoref{fig:hp-inter}所示。

\begin{figure}[H]
    \centering
    \includegraphics[width=\linewidth]{hp_interratio.pdf}
    \caption{\label{fig:hp-inter}\texttt{inter\_ratio} 在 $\varepsilon = 2.0$ 下的敏感性扫描}
\end{figure}

\par 由\autoref{fig:hp-inter}可见，扫描结果在四项指标上呈现明确的``trade-off''规律：
\textbf{NMI 与 Modularity Relative 随 \texttt{inter\_ratio} 增大而恶化}
（NMI 在 $[0.05, 0.20]$ 区段单调下降并在 $0.30$ 处趋于平稳，
Modularity Relative 单调上升），而 \textbf{CC Relative 则随之单调改善}；
Degree KL 在扫描范围内基本平坦。具体地：

\par 在 NMI 上，$0.05$（$0.336 \pm 0.006$，下文括号内均为 $95\%$ CI）是扫描范围内的最优取值，
显著优于其余四个取值（与 $0.10$、$0.15$、$0.20$、$0.30$ 的 $p$ 值分别为 $0.002$、$<0.001$、$<0.001$、$<0.001$）；
默认值 $0.10$（$0.319 \pm 0.008$）次之，并显著优于 $0.15$、$0.20$、$0.30$
（$p = 0.009$、$0.001$、$0.017$）。
在 Modularity Relative 上，$0.05$（$0.072 \pm 0.028$）同样是最优取值，
显著优于 $0.15$、$0.20$、$0.30$（$p = 0.003$、$<0.001$、$<0.001$），
但与默认值 $0.10$（$0.102 \pm 0.025$）之间无显著差异（$p = 0.116$）；
$0.10$ 同时显著优于 $0.20$、$0.30$（$p < 0.001$）。
在 CC Relative 上规律相反，$0.30$（$0.059 \pm 0.014$）显著优于所有其他取值，
$0.05$（$0.262 \pm 0.038$）则显著最差；
默认值 $0.10$（$0.180 \pm 0.036$）位于中间，显著优于 $0.05$（$p = 0.003$），
但显著劣于 $0.20$、$0.30$（$p = 0.004$、$<0.001$）。
在 Degree KL 上，扫描范围内绝大多数两两差异均不显著（$p > 0.05$），仅 $0.30$ 略优于 $0.05$、$0.20$。

\par 综合来看，扫描结果呈现典型的``社区结构 vs 局部传递性''的折中：
偏小的 \texttt{inter\_ratio} 通过更激进地缩减跨社区边来强化社区边界，
从而提升 NMI 与 Modularity Relative；但同时也使 Chung-Lu 采样的度序列放大幅度更大，
社区内部边密度上升，与原图聚类系数的偏离增大。
反之，偏大的 \texttt{inter\_ratio} 保留较多跨社区边、度序列放大幅度较小，
CC Relative 更接近原图，代价是 NMI 与 Modularity Relative 退化。
默认值 \texttt{inter\_ratio} = $0.10$ 是该折中曲线上的折中点：
在最关键的社区结构指标 NMI 上仅次于 $0.05$（差距 $0.017$，统计显著但绝对值较小），
在 Modularity Relative 上与最优值 $0.05$ 无显著差异，
而在 CC Relative 上保持中等水平。
若实际应用更看重社区结构指标，本节扫描表明 \texttt{inter\_ratio} = $0.05$ 是更具竞争力的选择；
若更看重局部传递性，则可考虑 $0.20$ 或 $0.30$。
本文初版调优时以 NMI 与 Modularity Relative 联合稳定为主要目标，
取 $0.10$ 作为默认值并完成所有主对比实验；将默认值进一步推向 $0.05$ 在
社区结构相关指标上仍有提升空间，留作后续工作。
\subsection{实验小结}
\label{subsec:exp-summary}

\par 本章通过主对比实验和超参敏感性分析，对本文提出的图重建优化方法
进行了系统评估。综合所有实验结果，可以得出以下两点结论：

\begin{enumerate}
    \item \textbf{整体可用性提升且具备一定可推广性。}
    在 Chameleon 数据集上，Ours 在 NMI 上将 7 个 $\varepsilon$ 的平均值
    由 PrivGraph 的 0.160 提升至 0.309（相对提升约 93\%），
    Modularity Relative 由 0.395 降至 0.121（相对降低约 69\%），
    Degree KL 由 4.12 降至 1.56（相对降低约 62\%），
    Diameter Relative 由 0.518 降至 0.390（相对降低约 25\%），
    在四项核心指标上稳定领先；
    EVC 类指标整体保持中性，差异在随机方差量级内。
    在 Facebook 上，本文方法在 NMI、Modularity Relative、
    Degree KL、CC Relative 四项指标上均稳定优于 PrivGraph，
    其中 Degree KL 平均由 3.98 降至 0.63、CC Relative 平均由 0.622 降至 0.301，
    改善尤为显著；
    CA-HepPh 上 NMI 平均由 0.066 提升至 0.376（5 倍以上），
    Modularity Relative 由 0.384 降至 0.157（相对降低约 59\%），
    Degree KL 与 CC Relative 同样保持全段领先。
    CA-HepPh 上 NMI 平均由 0.066 提升至 0.376（5 倍以上），
    Modularity Relative 由 0.384 降至 0.157（相对降低约 59\%），
    Degree KL 与 CC Relative 同样保持全段领先。
    在规模最大的 Enron 上，本文方法在社区结构相关指标上的改善最为显著——
    NMI 均值由 0.017 提升至 0.341（相对提升超过一个量级），
    Modularity Relative 由 0.435 降至 0.297（相对降低约 32\%）；
    但受 Enron 显著的 hub-and-spoke 拓扑特征影响，
    CC Relative 在该数据集上出现反常偏高，
    相关原因与缓解思路已在 \autoref{subsec:exp-main-other} 中给出分析。
    综合来看，本文方法的核心优势——
    基于度信息的跨社区采样与边预算补偿——
    在不同规模与不同来源图上稳定可推广。

    \item \textbf{核心超参的取值存在``社区结构 vs 局部传递性''的折中。}
    对 \texttt{inter\_ratio} 在 $\{0.05, 0.10, 0.15, 0.20, 0.30\}$ 上独立运行 30 次
    并做 Welch t 检验，扫描结果显示四项指标在取值方向上呈现一致的折中规律：
    NMI 与 Modularity Relative 倾向于偏小的取值（$0.05$ 是这两项的最优），
    而 CC Relative 倾向于偏大的取值（$0.30$ 最优），Degree KL 在范围内基本平坦。
    默认值 $0.10$ 处于折中曲线的中间区域：在 NMI 上仅次于 $0.05$（差距 $0.017$，$p = 0.002$）、
    在 Modularity Relative 上与最优值 $0.05$ 无显著差异（$p = 0.116$）、
    在 CC Relative 上居中，是兼顾多个指标的稳健选择。
    若实际应用更看重社区结构相关指标，扫描表明 $0.05$ 是更具竞争力的取值；
    若更看重局部传递性，则可考虑 $0.20$ 或 $0.30$。
\end{enumerate}

\par 与此同时，实验中也观察到一些尚未完全解决的问题。首先，
在 Chameleon 数据集上，CC Relative 这一局部传递性指标上 Ours 平均值 0.279
略高于 PrivGraph 的 0.206，
反映出度序列放大策略在补偿总边数的同时会使社区内部边密度局部上升、
合成图的聚类系数与原图的偏离轻微增大。
由于 Chung-Lu 模型本身缺乏对三角结构的显式建模，
这一问题在合成图层面难以彻底解决，本文将其作为局部三角结构修复的开放方向。
该现象在 Enron 数据集上被进一步放大——
Enron 具有显著的 hub-and-spoke 拓扑特征，
本文方法的度加权采样与度序列等比例放大共同强化了以 hub 为中心的星形子结构，
使 CC Relative 在 7 个 $\varepsilon$ 下均值达 2.93，远高于 PrivGraph 的 0.44。
初步实验表明，将三阶段隐私预算从默认均等分配调整为偏重信息抽取阶段的
$(0.25, 0.50, 0.25)$，可在多个预算下稳定降低 CC Relative 30\% 以上，
且不损害社区结构相关指标，
说明该副作用在算法层面尚可通过预算重配进行缓解。
其次，在极低预算（$\varepsilon = 0.5$）下 $\hat{v}$ 与 $\mathbb{C}_F$ 双重噪声过大，
合成图的结构指标整体偏差仍然较大，
本文方法虽较 PrivGraph 有改善，但绝对水平仍有提升空间。再次，
本文的核心超参 \texttt{inter\_ratio} 在 Chameleon 上进行了网格调优，
且本章敏感性扫描显示存在``社区结构 vs 局部传递性''的折中——
$0.05$ 在 NMI 与 Modularity Relative 上更具竞争力，
而本文取 $0.10$ 作为整体均衡的折中选择，
若将默认值进一步推向 $0.05$ 并在其他数据集上重新调优，
社区结构相关指标仍有进一步的提升空间。

\clearpage

\section{总结与展望}

\par 本章对全文工作进行总结，归纳本文的主要研究内容与贡献，
并对未来可能的研究方向进行展望。

\subsection{工作总结}

\par 本文以差分隐私图发布作为研究背景，
针对当前该方向中具有代表性的 PrivGraph 算法在图重建阶段所存在的局限性，
提出了一种基于度加权采样与边预算补偿的优化方法，
在不引入额外隐私开销的前提下提升了合成图的可用性。
全文的主要研究工作可以归纳为如下三个方面。

\par \textbf{（1）系统梳理了差分隐私图发布领域的研究现状。}
本文从图统计发布、图合成发布以及图数据合成的其他相关研究三条主线，
回顾了边差分隐私与节点差分隐私下的代表性工作，
对邻接矩阵扰动类（TmF）、密度划分类（DER）、概率图模型类（PrivHRG）、
本地差分隐私下的图合成方法（LDPGen、Block-HRG）以及
基于社区结构的方法（PrivGraph、PGB）进行了对比分析，
明确了 PrivGraph 在已有方法中的位置及其图重建阶段的可改进空间，
为本文工作的开展提供了清晰的研究脉络。

\par \textbf{（2）系统分析了 PrivGraph 图重建阶段的局限性。}
本文从噪声社区划分与噪声边向量两层误差耦合的角度，
对 PrivGraph 原算法在跨社区边处理上的问题进行了剖析：
原算法按噪声边向量 $\hat{v}$ 全量进行均匀采样的做法，
既未利用噪声度序列 $\hat{d}$ 中蕴含的节点重要性信息，
又会与已经被差分隐私机制扰动过的社区划分发生叠加效应；
跨社区边在 Louvain 等模块度优化算法中具有``桥接''作用，
忠实生成所有跨社区边会使合成图在重新执行社区检测时
与原图划分差异较大，进而拉低 NMI、模块度等社区结构相关指标。
该分析从机制层面揭示了原算法在低预算下合成图可用性下降的成因，
为后续优化设计明确了切入点。

\par \textbf{（3）提出并验证了面向跨社区边处理的图重建优化方法。}
针对上述局限性，本文对图重建阶段做出了两项相互配合的改动：
将原算法在跨社区方向上的均匀随机采样替换为基于噪声度序列 $\hat{d}$
的度加权采样，并按 \texttt{inter\_ratio}（默认 $0.10$）的比例
主动缩减每对社区之间生成的跨社区边数；
同时，在 Chung-Lu 采样之前根据每个社区因跨社区缩减所节省的边数
对 $\hat{d}$ 进行等比例放大，
使 Chung-Lu 模型在社区内部自然多生成相应数量的边，
在保持节点度分布形状不变的前提下维持总边数的稳定。
由于所有改动仅作用于已发布的噪声中间结果、不再访问原始数据，
根据差分隐私的后处理性质，
该方法不消耗额外的隐私预算，
整体隐私性与原始 PrivGraph 完全一致。
本文在 Chameleon、Facebook、CA-HepPh、Enron 四个公开数据集上
对所提方法进行了系统的实验评估，
结果表明该方法在 NMI、模块度相对误差等社区结构核心指标上
相对原始 PrivGraph 取得了稳定改善，
在不同规模、不同来源、不同度分布特性的图上展现出较好的可推广性。

\subsection{未来工作展望}

\par 本文工作虽然在合成图可用性上取得了一定改善，
但差分隐私图发布作为一个活跃的研究方向，
仍有许多值得进一步探索的问题。
结合本文工作的延伸方向，未来的研究可以从以下几个方面继续推进。

\par \textbf{（1）面向局部三角结构的合成图后处理修复。}
Chung-Lu 模型仅利用节点度信息生成边，
不显式建模节点之间的三角闭合关系，
使合成图在聚类系数等局部传递性指标上与原图之间存在系统性偏差。
未来可在保持差分隐私后处理性质的前提下，
引入三角闭合（triangle closure）等局部结构修复机制：
在合成图基础上以某种受控方式将开放的二阶路径 $u\text{-}v\text{-}w$
（即 $u \sim v$ 且 $v \sim w$ 但 $u \not\sim w$）
按一定概率闭合为三角形，从而提升合成图的聚类系数。
该类操作仅作用于已合成的图，符合差分隐私后处理性质，
不会引入额外隐私开销，是值得进一步探索的方向。

\par \textbf{（2）核心超参的数据自适应选取策略。}
本文的超参敏感性分析表明，
\texttt{inter\_ratio} 在不同取值下呈现出明确的
``社区结构 vs 局部传递性''折中规律，
而本文取统一默认值 $0.10$ 在不同拓扑特性的图上未必都是最优折中点。
未来可基于图的可公开估计特征
（如节点数、加噪后的平均度、加噪后的社区数与社区规模分布等）
建立 \texttt{inter\_ratio} 的经验选取规则或回归模型，
使方法在不同规模、不同拓扑特性的图上
均能自动取得接近最优的折中点；
进一步地，还可探索面向不同评测指标偏好的多目标超参选取策略，
使用户能够根据实际任务需求
灵活权衡社区结构与局部传递性两类目标。

\par \textbf{（3）三阶段隐私预算的联合最优分配。}
本文与 PrivGraph 一致，
将整体隐私预算 $\varepsilon$ 平均分配给
社区初始化、社区调整和信息抽取三个阶段。
然而对 Enron 数据集的初步实验表明，
适度偏重信息抽取阶段的预算分配可在保持社区结构指标既有优势的同时
抑制聚类系数指标上的副作用。
未来可对三阶段隐私预算的分配比例进行系统研究，
建立从图特征到最优分配的映射关系，
特别是针对具有显著 hub 结构、高聚类特性的网络
设计自适应预算分配机制。
该研究可与上一条方向结合，
形成``超参 + 预算''联合自适应的优化框架。

\par \textbf{（4）向更严格的隐私模型与更复杂的图形态推广。}
本文工作仅在中心化的边差分隐私模型下展开讨论。
将本文的核心思路——基于噪声中间结果的度加权采样与边预算补偿——
推广至本地差分隐私（LDP）、节点差分隐私等更严格的隐私模型，
或扩展到带属性图、动态图、异构图等更复杂的图形态，
是值得进一步探索的方向。
特别是在 LDP 模型下，由于每个用户仅能观测自身视图，
如何在分布式约束下复现本文方法的两项核心机制
将是后续研究的重要课题。

% \subsection{探索性尝试与失败分析}
% \label{subsec:exploratory}

% \par 在确定图重建作为主线优化方向之前，本文还围绕 PrivGraph 的不同阶段
% 进行了若干探索性尝试。这些尝试在初始动机上均有合理性，
% 但在系统的实验评估中未能展现出稳定、可推广的增益。
% 本节对其设计思路、实验观察与失败成因进行简要讨论，
% 为后续研究提供已被实证检验的参考路径。

% \subsubsection{基于零模型的跨社区边清理后处理}

% \par \textbf{设计动机：}\autoref{chap:limitation} 中指出，PrivGraph 合成图存在
% 跨社区边过多的现象，且其中相当比例落在度数较低、本不应连通的节点对之间。
% 受配置零模型（configuration null model）的启发，本文设计了一种合成图后处理：
% 对每条跨社区边 $(u, v)$，计算其在零模型下的期望连接数
% $\mathrm{expected}(u, v) = \hat{d}_u \hat{d}_v / (2m)$；
% 对 $\mathrm{expected}(u, v) \leq 1$ 的边判定为``虚假桥梁''并以一定比例
% \texttt{swap\_ratio} 进行重排或删除，期望以此清理稀疏跨社区边、
% 间接修复局部三角结构。

% \par \textbf{实验观察：}在 Chameleon 数据集上，单独启用该模块
% 相对 PrivGraph 在 Modularity Relative（平均由 0.408 降至 0.294，相对降低约 28\%）、
% CC Relative（平均由 0.227 降至 0.165，相对降低约 27\%）以及
% Degree KL（平均由 4.17 降至 1.65，相对降低约 60\%）上能取得局部改善；
% 但在最关键的社区结构指标 NMI 上几乎无变化（0.155 vs 0.158），
% 且与图重建模块组合后在多数指标上反而引入小幅扰动。

% \par \textbf{超参敏感性扫描：}\autoref{fig:hp-swap} 给出了 \texttt{swap\_ratio}
% 在 $\{0.0, 0.1, 0.3, 0.5, 0.7\}$ 上 30 次独立重复的扫描结果。
% NMI 上 $0.1$（$0.318 \pm 0.007$，下文括号内均为 $95\%$ CI）显著低于
% $0.5$（$0.332 \pm 0.007$，$p = 0.007$）和 $0.7$（$0.335 \pm 0.008$，$p = 0.003$），
% 但默认 $0.3$（$0.324 \pm 0.009$）与上述各取值之间均无显著差异；
% Modularity Relative 上 $0.0$（$0.135 \pm 0.036$）显著优于
% $0.1$（$0.216 \pm 0.047$，$p = 0.011$），$0.5$（$0.146 \pm 0.047$）也显著优于
% $0.1$（$p = 0.043$），但默认 $0.3$（$0.180 \pm 0.056$）与其他取值之间均无显著差异；
% Degree KL 与 CC Relative 上各两两差异均不显著。
% 有意义的发现是 $\texttt{swap\_ratio} = 0.1$ 在 NMI 与 Modularity Relative 上
% 一致地显著劣于多个其他取值，说明``少量''边交换反而会引入度扰动而无明显收益；
% 扫描范围内未能找到全指标一致占优的设置。

% \begin{figure}[!t]
%     \centering
%     \includegraphics[width=\linewidth]{hp_swapratio.pdf}
%     \caption{\label{fig:hp-swap}\texttt{swap\_ratio} 在 $\varepsilon = 2.0$ 下的敏感性扫描}
% \end{figure}

% \par \textbf{失败分析：}进一步分析发现两个问题：
% 其一，``$\mathrm{expected} \leq 1$''阈值依赖于 $\hat{d}$ 与 $\hat{v}$，
% 在低预算下两者噪声较大，使``虚假桥梁''的判定本身已不够可靠；
% 其二，删除/重排操作对节点度的局部扰动会反向作用于已通过 Chung-Lu
% 和度序列放大精心构造的度分布，使两个模块的目标不一致。
% 综合来看，该后处理思路与图重建模块在作用对象上有一定耦合，
% 难以稳定地``纯叠加''，故本文未将其纳入主线方案。

% \subsubsection{基于 DP 度+邻居得分的社区初始化}

% \par \textbf{设计动机：}原 PrivGraph 在社区初始化阶段将所有节点随机分组，
% 每组 $N$ 个节点；该步骤完全不利用节点结构信息，
% 直觉上若能以差分隐私方式利用节点的度信息进行非随机分组，
% 后续 Louvain 阶段或可获得更好的起点。
% 本文据此设计了基于 DP 度 + 一阶邻居得分的初始化算法：
% 将隐私预算 $\varepsilon_1$ 划分为得分预算与初始化预算两部分，
% 通过加噪度 $\hat{d}_u$ 与一阶邻居得分 $\widehat{ns}_u = \sum_{v \in N(u)} \hat{d}_v$
% 的加权组合 $\hat{d}_u + \beta \cdot \widehat{ns}_u$ 对节点排序，
% 按得分顺序依次分组，期望使得分相近的节点（即结构相似的节点）落在同一组中。

% \par \textbf{实验观察与失败分析：}
% 在 Chameleon 上的初步对比表明，该方法相对原始随机分组在 NMI 上没有稳定改善，
% 部分预算下反而轻微下降。失败的主要原因有两点：
% 其一，邻居得分的 $\ell_1$ 敏感度为 $2C$（$C$ 为度的截断阈值），
% 要保证敏感度有界必须对度做较紧的截断，而截断本身会损失高度数节点的信息；
% 其二，得分排序仅利用了局部度信息，未利用拓扑结构，
% 对于社区结构鲜明的图，按度排序得到的分组与真实社区划分相关性较弱，
% 反而在 Louvain 起点上引入了系统性偏差。基于此，
% 本文未将该模块纳入主线。

% \subsubsection{基于经验公式的隐私预算动态分配}

% \par \textbf{设计动机：}原 PrivGraph 将整体预算 $\varepsilon$
% 平均分配给社区初始化、社区调整、信息抽取三个阶段（$1{:}1{:}1$）。
% 直觉上，三个阶段对噪声的敏感程度并不相同——
% 信息抽取阶段直接决定 $\hat{d}$ 与 $\hat{v}$ 的精度，
% 对图重建质量影响最大；社区初始化的噪声又可以在社区调整阶段部分修复。
% 此外，最优分配很可能依赖于具体的 $(\varepsilon, n, d, K)$ 配置，
% 而非一个固定常数比。本文据此尝试以数据驱动的方式探索更优的预算分配方案。

% \par \textbf{探索方法：}本文采取``热力图扫描 + 经验公式拟合''的两阶段流程。
% 首先在 Chameleon 上对 $(\varepsilon_1, \varepsilon_2, \varepsilon_3)$ 三元组
% （约束 $\sum \varepsilon_i = \varepsilon$）进行细粒度网格扫描，
% 在多个 $\varepsilon$ 下绘制 NMI、Modularity Relative、Degree KL 等核心指标
% 随分配比例变化的热力图，观察是否存在系统性的``最优区域''。
% 其次，以扫描得到的较优区域为目标，
% 以 $\varepsilon$、节点数 $n$、平均度 $d$、社区数 $K$、
% 以及由 $\hat{d}$ 与 $\varepsilon_3$ 估计的社区内信噪比 SNR 为特征，
% 拟合一组经验公式 $(e_1^*, e_2^*, e_3^*)$，
% 并通过截断与归一化保证合法性。
% 为利用社区划分完成后才可观测的 $K$ 与 SNR 信息，
% 该公式被进一步组织为两阶段：
% 首阶段基于 $(\varepsilon, n, d)$ 给出 $\varepsilon_1$ 与初步的 $(\varepsilon_2, \varepsilon_3)$，
% 社区初始化结束后再根据 $K$ 与社区结构的估计动态修正
% $\varepsilon_2 : \varepsilon_3$ 的剩余分配。

% \par \textbf{实验观察与失败分析：}
% 该方案在 Chameleon 上的部分 $\varepsilon$ 取值上相对 $1{:}1{:}1$ 默认分配
% 确实取得了轻微改善，但综合多数据集与多指标评估后暴露出两个根本性问题。

% \par \textit{其一，跨数据集不可推广。}
% 基于 Chameleon 拟合的经验公式在迁移到 Facebook、CA-HepPh 等数据集后表现不稳定，
% 部分预算下甚至相对默认分配轻微退化。
% 进一步检查热力图发现，``较优区域''在不同数据集上的位置本身就有显著差异——
% 不同的拓扑结构（社区规模分布、度分布尾部、聚类系数）使得三阶段噪声
% 对最终合成图质量的``边际贡献''比例并不一致，
% 因此在单一数据集上拟合的公式难以泛化。
% 若要为每个数据集单独调参，则违背了``数据集无关的隐私预算策略''这一初衷。

% \par \textit{其二，不存在多指标一致占优的分配点。}
% 即使在 Chameleon 内部，热力图也显示出明显的``指标互相冲突''现象：
% 有利于 NMI、Modularity Relative 的分配区域通常使 Degree KL、CC Relative 退化，
% 反之亦然。这与本文图重建模块超参敏感性扫描所揭示的
% ``社区结构 vs 局部传递性''折中具有相似的内在结构——
% 将更多预算分配给信息抽取会更准确地恢复 $\hat{d}$ 与 $\hat{v}$
% （更利于度分布与局部结构指标），但相应削弱社区划分的精度
% （不利于 NMI 与模块度）。因此在 Pareto 意义上不存在单一``最优''分配点，
% 只存在一条贯穿多个指标的折中前沿。

% \par 综合以上两点，``为任意 $(\varepsilon, n, d, K)$ 拟合一个通用的最优分配公式''
% 在 PrivGraph 现有框架下难以实现。
% 更准确的结论是：预算分配在多指标空间中沿一条 Pareto 前沿滑动，
% $1{:}1{:}1$ 默认分配恰好落在该前沿上较为中性的位置，
% 是简单且较均衡的稳健选择。
% 若实际应用对某类指标有明确偏好，
% 可基于本文的扫描思路在目标数据集上做局部微调；
% 但本文未能找到对所有数据集与所有指标一致占优的公式，
% 故未将该模块纳入主线方案。